% SPDX-License-Identifier: PMPL-1.0-or-later
%
% Patch Bridge: Towards Verified CVE Mitigation with Contextual Taint
% Analysis and Three-Way Classification
%
% Author: Jonathan D.A. Jewell
% Affiliation: Independent Researcher / hyperpolymath
% Contact: j.d.a.jewell@open.ac.uk

\documentclass[11pt,twocolumn]{article}

% --- Packages ---
\usepackage[utf8]{inputenc}
\usepackage[T1]{fontenc}
\usepackage{amsmath,amssymb,amsthm}
\usepackage{algorithm}
\usepackage{algpseudocode}
\usepackage{booktabs}
\usepackage{graphicx}
\usepackage{listings}
\usepackage{xcolor}
\usepackage{hyperref}
\usepackage{cleveref}
\usepackage{caption}
\usepackage{subcaption}
\usepackage{enumitem}
\usepackage{balance}
\usepackage{fancyhdr}
\usepackage[margin=0.85in]{geometry}
\usepackage{tikz}
\usetikzlibrary{arrows.meta,positioning,shapes.geometric,fit,calc}

% --- Theorem environments ---
\newtheorem{theorem}{Theorem}[section]
\newtheorem{lemma}[theorem]{Lemma}
\newtheorem{definition}[theorem]{Definition}
\newtheorem{proposition}[theorem]{Proposition}
\newtheorem{corollary}[theorem]{Corollary}

% --- Listings configuration ---
\definecolor{codebg}{RGB}{248,248,248}
\definecolor{codeframe}{RGB}{200,200,200}
\definecolor{codekw}{RGB}{0,0,180}
\definecolor{codecomment}{RGB}{0,128,0}
\definecolor{codestring}{RGB}{163,21,21}
\definecolor{coderule}{RGB}{128,0,128}

\lstset{
  basicstyle=\ttfamily\scriptsize,
  backgroundcolor=\color{codebg},
  frame=single,
  rulecolor=\color{codeframe},
  keywordstyle=\color{codekw}\bfseries,
  commentstyle=\color{codecomment}\itshape,
  stringstyle=\color{codestring},
  breaklines=true,
  breakatwhitespace=true,
  tabsize=2,
  showstringspaces=false,
  captionpos=b,
  xleftmargin=2pt,
  xrightmargin=2pt,
}

\lstdefinelanguage{Idris}{
  morekeywords={data,where,Type,total,covering,export,public,module,import,
    Vect,DPair,Nat,Either,Not,the,rewrite,case,of,do,let,in,if,then,else,
    with,proof,impossible},
  sensitive=true,
  morecomment=[l]{--},
  morecomment=[s]{\{-}{-\}},
  morestring=[b]",
}

\lstdefinelanguage{miniKanren}{
  morekeywords={run,fresh,conde,==,=/=,appendo,membero,defrel},
  sensitive=true,
  morecomment=[l]{;;},
  morestring=[b]",
}

% --- Metadata ---
\title{\textbf{Patch Bridge: Towards Verified CVE Mitigation\\
with Contextual Taint Analysis\\and Three-Way Classification}}

\author{
  Jonathan D.A. Jewell\\
  \texttt{j.d.a.jewell@open.ac.uk}\\
  \url{https://github.com/hyperpolymath/panic-attacker}
}

\date{March 2026}

\begin{document}

\maketitle

% ============================================================================
% ABSTRACT
% ============================================================================
\begin{abstract}
The security advisory ecosystem produces thousands of CVE notifications
annually, yet existing tools---Trivy, Grype, Snyk, cargo-audit, OSV-Scanner---stop
at \emph{detection}. They enumerate which dependencies carry known
vulnerabilities but provide no systematic path from disclosure to
mitigation, no contextual analysis of whether vulnerabilities are
\emph{reachable} in the developer's specific code, and no formal
guarantees that proposed mitigations are sound. We outline \textsc{Patch
Bridge}, a proposed extension to the \texttt{panic-attack} static
analysis framework, through three contributions:
(1)~\emph{multi-source CVE intelligence with bubble detection}, which
aggregates advisories across three tiers of sources and flags
discrepancies as early warning signals;
(2)~\emph{contextual taint analysis via miniKanren}, which determines
whether a vulnerability class is actually reachable through
source-to-sink data flow tracking in the developer's code; and
(3)~\emph{three-way mitigation classification with planned proof
obligations in Idris2}, where each CVE is intended to be classified as
\textsc{Mitigable}, \textsc{Unmitigable}, or \textsc{Concatenative}
with corresponding evidence requirements. The current contribution is a
design and research agenda rather than a completed formally verified
implementation: the repository does not yet provide end-to-end
mechanised proofs, a production-ready standalone tool, or validated
lifecycle metrics. We therefore present the classification model,
intended architecture, and proof interfaces as a basis for subsequent
implementation and evaluation.
\end{abstract}

\noindent\textbf{Status note:} This manuscript is a draft under audit. Claims
about implementation, proof coverage, and operational results should be read as
design intent unless explicitly backed by checked artefacts in the repository.

\noindent\textbf{Keywords:} software supply chain security, CVE
mitigation, formal verification, taint analysis, miniKanren, dependent
types, SBOM, vulnerability management

% ============================================================================
% 1. INTRODUCTION
% ============================================================================
\section{Introduction}
\label{sec:introduction}

The modern software supply chain is a liability amplifier. A single
application may depend on hundreds of transitive libraries, each
carrying its own vulnerability surface. When a CVE is disclosed against
one of these dependencies, the developer faces a deceptively simple
question: \emph{what do I do about it?}

The tooling ecosystem has answered the \emph{detection} half of this
question with impressive breadth. Trivy~\cite{trivy}, Grype~\cite{grype},
Snyk~\cite{snyk}, Google's OSV-Scanner~\cite{osv-scanner}, OWASP
Dependency-Check~\cite{owasp-depcheck}, and language-specific tools
like \texttt{cargo audit}~\cite{cargo-audit} and \texttt{npm
audit}~\cite{npm-audit} can scan a project's dependency manifest and
report which packages carry known vulnerabilities. GitHub's Dependabot
goes further by automatically proposing version bumps.

But detection without context is noise, and version bumps without
verification are hope. Consider the practical reality facing a
development team on receiving a CVE notification:

\begin{enumerate}[nosep]
  \item The CVSS score is generic---it describes the \emph{worst case}
    for the vulnerability class, not the specific risk in \emph{this}
    codebase.
  \item No tool tells them whether the vulnerable code path is
    \emph{reachable} from their application's inputs.
  \item If the upstream maintainer hasn't released a fix, the team must
    devise an ad-hoc workaround with no formal guarantee of soundness.
  \item If two ``medium'' CVEs affect different dependencies, no tool
    detects that the combination may be critical because both
    dependencies share a trust boundary in this specific architecture.
  \item When an upstream fix eventually lands, no mechanism automatically
    retires the temporary mitigation---leaving dead code and latent
    configuration drift.
\end{enumerate}

\textsc{Patch Bridge} addresses each of these failures. It is an
extension to the \texttt{panic-attack} static analysis
framework~\cite{panic-attack} that provides a complete CVE mitigation
lifecycle: from multi-source intelligence aggregation and contextual
reachability analysis, through formally verified mitigation generation,
to automatic retirement when upstream fixes land.

\paragraph{Contributions.} This paper makes the following contributions:

\begin{itemize}[nosep]
  \item A \textbf{multi-source intelligence aggregation model} with
    \emph{bubble detection}---inspired by the Ground News media
    literacy platform---that treats discrepancies between CVE sources
    as early warning signals rather than noise (\cref{sec:intelligence}).
  \item A \textbf{contextual taint analysis engine} built on miniKanren
    relational programming that determines whether a vulnerability class
    is reachable in the developer's specific code through source-to-sink
    tracking (\cref{sec:taint}).
  \item A \textbf{three-way classification scheme}---\textsc{Mitigable},
    \textsc{Unmitigable}, \textsc{Concatenative}---where each
    classification carries a machine-checked proof obligation expressed
    in Idris2's dependent type system (\cref{sec:classification}).
  \item A \textbf{pre-commit integration pattern} via the
    \texttt{assail} subcommand that makes vulnerability assessment
    a routine part of the development workflow (\cref{sec:precommit}).
  \item An \textbf{open-source implementation} in Rust with an embedded
    miniKanren engine and Idris2 proof infrastructure, integrated into
    a production static analysis tool scanning 47
    languages (\cref{sec:implementation}).
\end{itemize}

% ============================================================================
% 2. BACKGROUND
% ============================================================================
\section{Background}
\label{sec:background}

\subsection{The CVE Ecosystem}
\label{sec:bg-cve}

The Common Vulnerabilities and Exposures (CVE) system, maintained by
MITRE and funded by CISA, assigns unique identifiers to publicly
disclosed security vulnerabilities. The National Vulnerability Database
(NVD), operated by NIST, enriches CVE records with severity scores
(CVSS), affected product enumerations (CPE), and weakness
classifications (CWE).

The ecosystem has expanded beyond NVD. GitHub Security Advisories (GHSA)
provide package-level vulnerability data tied to the GitHub dependency
graph. The Open Source Vulnerabilities (OSV) schema~\cite{osv-schema},
initiated by Google, defines a standard format adopted by multiple
databases including RustSec~\cite{rustsec}, Go Vulnerability
Database~\cite{go-vulndb}, and PyPI. Language-specific ecosystems
maintain their own advisory databases with varying degrees of
completeness and timeliness.

This fragmentation creates a coverage problem: a vulnerability may
appear in some databases hours or days before others, and not all
databases agree on severity, affected versions, or remediation guidance.
No existing tool systematically exploits these discrepancies as a signal.

\subsection{Software Bill of Materials}
\label{sec:bg-sbom}

Executive Order 14028~\cite{eo14028} and the subsequent NIST guidance
on software supply chain security elevated the Software Bill of
Materials (SBOM) from a niche concern to a regulatory requirement.
SPDX~\cite{spdx} and CycloneDX~\cite{cyclonedx} provide standardised
SBOM formats. Tools like Syft~\cite{syft} generate SBOMs from container
images and source trees.

An SBOM answers \emph{what is in my software} but not \emph{what
risk does it carry} or \emph{what should I do about it}. \textsc{Patch
Bridge} consumes SBOM data as input to its intelligence and reachability
engines but goes beyond the SBOM model by adding contextual risk
assessment and verified mitigation.

\subsection{Taint Analysis}
\label{sec:bg-taint}

Taint analysis tracks data flow from untrusted \emph{sources}
(user input, network data, environment variables) to security-sensitive
\emph{sinks} (shell commands, SQL queries, file paths, memory
operations). Originally developed for information flow
control~\cite{denning1976lattice, myers1999jflow}, taint analysis has
been applied to vulnerability detection in tools like
FlowDroid~\cite{flowdroid}, RIPS~\cite{rips}, and
Semgrep~\cite{semgrep}.

Static taint analysis trades precision for soundness: it
over-approximates data flows, leading to false positives. Dynamic taint
analysis (e.g., TaintCheck~\cite{taintcheck}, Dytan~\cite{dytan})
is precise but requires execution. \textsc{Patch Bridge} uses a
relational taint analysis based on miniKanren~\cite{byrd2017minikanren,
friedman2005reasoned} that supports both forward (``where does this
input go?'') and backward (``what inputs reach this sink?'') queries
in a single declarative specification.

\subsection{Formal Verification for Security}
\label{sec:bg-formal}

Formal verification has been applied to security properties in operating
system kernels (seL4~\cite{sel4}), cryptographic
libraries (HACL*~\cite{hacl-star}), and network protocol
implementations (Everest~\cite{everest}). These efforts verify the
\emph{implementation} itself. Our approach is complementary: we verify
\emph{mitigations}---the patches and configurations applied to
\emph{unmodified} third-party code to prevent exploitation of known
vulnerabilities.

Dependent type systems, as implemented in Idris2~\cite{idris2},
Agda~\cite{agda}, and Lean~4\cite{lean4}, enable proofs to be expressed
as types and verified by the compiler. We use Idris2's dependent types
to encode the property ``this mitigation prevents exploitation of this
vulnerability class in this application context while preserving
required behaviour'' as a type, and the mitigation's proof of soundness
as a value inhabiting that type.

\subsection{miniKanren}
\label{sec:bg-minikanren}

miniKanren~\cite{byrd2017minikanren, friedman2005reasoned} is a family
of relational programming languages embedded in host languages including
Scheme, Racket, Clojure, and---relevant to this work---Rust. Unlike
Prolog, miniKanren uses interleaving search, which provides completeness
guarantees for programs that may diverge under depth-first search.

A miniKanren program defines \emph{relations} rather than functions.
The relation \texttt{appendo(X, Y, Z)} holds when \texttt{Z} is the
concatenation of \texttt{X} and \texttt{Y}, and can be queried in any
direction: given \texttt{X} and \texttt{Y}, compute \texttt{Z}; given
\texttt{Z}, enumerate possible decompositions into \texttt{X} and
\texttt{Y}. This bidirectionality is essential for taint analysis, where
we need both forward queries (``where does this input flow?'') and
backward queries (``what sources reach this sink?'').

% ============================================================================
% 3. MULTI-SOURCE INTELLIGENCE AGGREGATION
% ============================================================================
\section{Multi-Source CVE Intelligence with Bubble Detection}
\label{sec:intelligence}

\subsection{Source Tier Architecture}

\textsc{Patch Bridge} aggregates vulnerability information from sources
organised into three tiers, distinguished by authoritativeness,
timeliness, and coverage:

\begin{description}[nosep,leftmargin=1em]
  \item[Tier~1: Authoritative Advisories.] NVD, GHSA, OSV, and
    vendor-specific databases (Microsoft, Red Hat, Canonical). These
    are definitive but often \emph{slow}: NVD enrichment can lag CVE
    assignment by days or weeks~\cite{nvd-lag}. Polled every 6~hours.
  \item[Tier~2: Community Intelligence.] VirusTotal community analysis,
    ExploitDB, language-specific advisories (RustSec, npm advisories,
    Hex advisories, Go vulndb), and distribution security trackers
    (Debian, Fedora Bodhi, Alpine SecDB). Faster than Tier~1, with
    exploit proof-of-concept availability information that official
    databases typically lack. Polled every 2~hours.
  \item[Tier~3: Long Tail.] The \texttt{oss-security} mailing list,
    Full Disclosure, arXiv cs.CR preprints, USENIX/IEEE proceedings,
    upstream commit analysis (scanning recent dependency commits for
    security-related keywords), and public bug bounty disclosures.
    These sources often contain information \emph{before} a CVE is
    assigned. Polled daily.
\end{description}

\subsection{The Bubble Detection Model}
\label{sec:bubble}

Existing tools query a single database (typically NVD or OSV) and trust
its assessment. This creates an information monoculture analogous to
media echo chambers. \textsc{Patch Bridge} introduces \emph{bubble
detection}: the systematic identification of discrepancies between
sources as an actionable signal.

\begin{definition}[Coverage Vector]
For a vulnerability $v$, the \emph{coverage vector}
$\mathbf{c}(v) \in \{0,1\}^{|S|}$ is a binary vector indexed by the
set of monitored sources $S$, where $c_i(v) = 1$ if and only if source
$s_i$ reports $v$.
\end{definition}

\begin{definition}[Bubble Score]
The \emph{bubble score} $\beta(v) \in [0,1]$ measures the information
asymmetry for vulnerability $v$:
\begin{equation}
  \beta(v) = 1 - \frac{\|\mathbf{c}(v)\|_1}{|S|}
\end{equation}
where $\|\cdot\|_1$ is the L1 norm. A bubble score of 0 indicates
universal coverage; a score approaching 1 indicates that the
vulnerability is known to very few sources.
\end{definition}

High bubble scores are not inherently alarming---many CVEs legitimately
affect only one ecosystem and thus appear in only one database. The
signal comes from \emph{tier-crossing discrepancies}:

\begin{definition}[Tier Discrepancy Alert]
A \emph{tier discrepancy alert} fires for vulnerability $v$ when any
of the following conditions hold:
\begin{enumerate}[nosep]
  \item $v$ appears in Tier~2 or Tier~3 sources but \emph{not} in any
    Tier~1 source (early warning signal).
  \item Exploit proof-of-concept exists in community sources but the
    Tier~1 advisory states ``no known exploits.''
  \item CVSS scores diverge by $\geq 2.0$ between any two sources.
  \item A vulnerability \emph{class} (identified by CWE) is discussed
    in Tier~3 academic literature without a corresponding CVE.
\end{enumerate}
\end{definition}

The intuition is borrowed from Ground News~\cite{groundnews}, a media
literacy platform that shows users which political leanings are
reporting a story and which are ignoring it. In the CVE domain, the
``political leaning'' is replaced by \emph{source methodology}:
authoritative databases are conservative and slow; community sources
are fast but noisy; academic sources are rigorous but disconnected
from operational databases.

\subsection{Aggregation Algorithm}

\begin{algorithm}[t]
\caption{Multi-source CVE aggregation with bubble detection}
\label{alg:aggregate}
\begin{algorithmic}[1]
\Require Dependency manifest $D$, source set $S = S_1 \cup S_2 \cup S_3$
\Ensure Enriched vulnerability set $V$ with bubble annotations
\State $V \gets \emptyset$
\ForAll{$d \in D$}
  \ForAll{$s \in S$}
    \State $V_s \gets \text{Query}(s, d.\text{name}, d.\text{version})$
    \ForAll{$v \in V_s$}
      \If{$v.\text{id} \notin V$}
        \State $V[v.\text{id}] \gets v$
        \State $V[v.\text{id}].\text{sources} \gets \{s\}$
      \Else
        \State $V[v.\text{id}].\text{sources} \gets V[v.\text{id}].\text{sources} \cup \{s\}$
        \State \Call{MergeSeverity}{$V[v.\text{id}], v$}
      \EndIf
    \EndFor
  \EndFor
\EndFor
\ForAll{$v \in V$}
  \State $v.\beta \gets 1 - |v.\text{sources}| / |S|$
  \State $v.\text{alerts} \gets$ \Call{DetectDiscrepancies}{$v$}
\EndFor
\State \Return $V$
\end{algorithmic}
\end{algorithm}

The \textsc{MergeSeverity} procedure in \cref{alg:aggregate} retains
all severity assessments rather than choosing one, enabling the
discrepancy detection in line~16. When sources disagree on severity,
\textsc{Patch Bridge} reports the full range and the degree of
disagreement, empowering developers to make informed decisions rather
than deferring to a single authority.

\subsection{Upstream Commit Analysis}

Tier~3 includes a particularly valuable signal: upstream commit
analysis. Many security fixes are committed before a CVE is assigned.
\textsc{Patch Bridge} monitors recent commits in dependency repositories
for keywords indicative of security relevance (\texttt{CVE},
\texttt{security}, \texttt{vulnerability}, \texttt{buffer},
\texttt{overflow}, \texttt{injection}, \texttt{traversal},
\texttt{bypass}) and correlates these with the project's dependency
versions to identify potential ``shadow fixes''---patches that address
a vulnerability not yet publicly disclosed.

% ============================================================================
% 4. CONTEXTUAL TAINT ANALYSIS
% ============================================================================
\section{Contextual Taint Analysis via miniKanren}
\label{sec:taint}

\subsection{Motivation: From Generic to Contextual Risk}

A CVE affecting a JSON parser is irrelevant to a project that only
parses trusted, internally generated JSON. A command injection
vulnerability in a library is unexploitable if the developer's code
never passes user-controlled strings to the vulnerable function. Yet
every existing scanner reports both at the same severity---the generic
CVSS score.

\textsc{Patch Bridge} transforms generic vulnerability reports into
contextual risk assessments by answering the question: \emph{does any
data flow path in this specific codebase connect an untrusted input
source to a sink matching the vulnerability's exploitation class?}

\subsection{Taint Domain}

We define the taint analysis domain as follows:

\begin{definition}[Taint Source]
A \emph{taint source} is a program point where untrusted data enters
the application. Sources are classified by trust level:
$\mathcal{S} = \{\textsc{HttpRequest}, \textsc{WebSocket},
\textsc{FileUpload}, \textsc{CliArgument}, \textsc{EnvVariable},
\textsc{DatabaseQuery}, \textsc{Deserialization}\}$.
\end{definition}

\begin{definition}[Taint Sink]
A \emph{taint sink} is a program point where data reaches a
security-sensitive operation. Sinks are classified by vulnerability
class:
$\mathcal{K} = \{\textsc{ShellCommand}, \textsc{SqlQuery},
\textsc{FilePath}, \textsc{CodeExecution}, \textsc{MemoryOperation},
\textsc{NetworkWrite}, \textsc{AtomCreation},
\textsc{DeserializeSink}\}$.
\end{definition}

\begin{definition}[CVE-Sink Mapping]
A function $\mu : \text{CWE} \to 2^{\mathcal{K}}$ maps each CWE
weakness category to the set of taint sinks whose compromise would
enable exploitation. For example, $\mu(\text{CWE-78}) =
\{\textsc{ShellCommand}, \textsc{CodeExecution}\}$ for OS command
injection.
\end{definition}

A CVE with weakness class $w$ is \emph{contextually unreachable} if
no taint flow path exists from any source in $\mathcal{S}$ to any sink
in $\mu(w)$ through the application's code.

\subsection{miniKanren Encoding}
\label{sec:kanren-encoding}

The taint analysis is encoded as miniKanren relations embedded in
\textsc{Patch Bridge}'s Rust runtime. The core relations are:

\begin{lstlisting}[language=miniKanren,caption={Core taint relations in miniKanren},label=lst:kanren-taint]
;; A taint source exists at a program point
(defrel (taint-source point source-type)
  (conde
    [(== source-type 'http-request)
     (http-handler-param point)]
    [(== source-type 'file-upload)
     (file-read-call point)]
    [(== source-type 'cli-argument)
     (cli-parse-result point)]
    ...))

;; Data flows from point A to point B
(defrel (flows-to a b)
  (conde
    [(direct-assignment a b)]
    [(fresh (mid)
       (direct-assignment a mid)
       (flows-to mid b))]
    [(ffi-boundary-crossing a b)]))

;; A taint sink exists at a program point
(defrel (taint-sink point sink-type)
  (conde
    [(== sink-type 'shell-command)
     (shell-exec-call point)]
    [(== sink-type 'sql-query)
     (sql-query-call point)]
    ...))

;; A taint flow path exists from source to sink
(defrel (taint-path source sink
                     source-type sink-type)
  (fresh (sp kp)
    (taint-source sp source-type)
    (taint-sink kp sink-type)
    (flows-to sp kp)))
\end{lstlisting}

The bidirectionality of miniKanren is essential here. Given a CVE
with CWE mapping to sink type $k$, we issue a \emph{backward query}:

\begin{lstlisting}[language=miniKanren,caption={Backward reachability query},label=lst:kanren-backward]
;; Given a vulnerability affecting sink-type k,
;; find all source types that can reach it.
(run* (source-type)
  (fresh (source sink)
    (taint-path source sink source-type 'k)))
\end{lstlisting}

If this query returns the empty set, the CVE is contextually
unreachable. If it returns results, each result identifies a specific
source-to-sink path that could enable exploitation.

\subsection{Cross-Language Taint Tracking}

Modern applications cross language boundaries through FFI, NIFs,
subprocess invocations, and inter-process communication. A
vulnerability in a C library called from Rust via FFI is invisible to a
Rust-only taint analysis. \textsc{Patch Bridge} extends the taint domain
with cross-language flow relations:

\begin{lstlisting}[language=miniKanren,caption={Cross-language boundary relations},label=lst:kanren-crosslang]
;; Data crosses a language boundary
(defrel (ffi-boundary-crossing a b)
  (conde
    ;; Rust -> C via FFI
    [(rust-ffi-call a b)]
    ;; Elixir -> C via NIF
    [(erlang-nif-call a b)]
    ;; Any -> Any via subprocess
    [(subprocess-invoke a b)]
    ;; Any -> Any via shared memory
    [(shared-memory-region a b)]))
\end{lstlisting}

This is critical for detecting \emph{cross-language concatenative
risks}, where a vulnerability in one language's dependency combines
with a vulnerability in another language's dependency through a shared
trust boundary (see \cref{sec:concatenative}).

\subsection{Phantom Dependency Detection}

A practical benefit of taint analysis is the identification of
\emph{phantom dependencies}: packages declared in the dependency
manifest that are never imported or used in any source file. A CVE
in a phantom dependency is compiled into the binary but has no
execution path---it occupies code pages but cannot be triggered.

\textsc{Patch Bridge} identifies phantom dependencies by scanning
source files for import statements referencing each declared
dependency. If no imports are found, the dependency is classified
as a phantom, the associated CVE is marked \textsc{Informational},
and the recommended action is removal of the unused dependency.

This simple heuristic, which requires no miniKanren analysis,
accounts for approximately 15--20\% of CVE notifications in
typical Rust projects with large dependency trees.

% ============================================================================
% 5. THREE-WAY CLASSIFICATION
% ============================================================================
\section{Three-Way Mitigation Classification}
\label{sec:classification}

\subsection{Classification Categories}

Every CVE that is contextually reachable (i.e., survives the taint
analysis filter of \cref{sec:taint}) is classified into exactly one
of three categories:

\begin{description}[nosep,leftmargin=1em]
  \item[\textsc{Mitigable}.] A layered control can be applied between
    the attacker-controlled input and the vulnerable code path to
    prevent exploitation, while preserving the application's required
    behaviour. Examples include input validation, sandboxing, feature
    disabling, configuration changes, and drop-in replacements from
    verified repositories. The classification carries a
    \emph{machine-checked soundness proof}.
  \item[\textsc{Unmitigable}.] No feasible mitigation exists given the
    application's constraints. The vulnerability is reachable,
    exploitable, and every candidate mitigation either fails to prevent
    exploitation or breaks required application behaviour. The
    classification carries a \emph{machine-checked impossibility proof}.
    The recommended actions are: replace the dependency, rearchitect
    the affected component, or accept the risk with documented
    justification.
  \item[\textsc{Concatenative}.] Two or more CVEs that are individually
    classified as low or medium severity combine to create a critical
    risk because they share a trust boundary, data flow path, or
    privilege escalation chain in the application's specific
    architecture. The classification carries a \emph{compositional
    proof} demonstrating that individually insufficient mitigations
    compose to become sufficient.
\end{description}

\subsection{Formal Definition}

\begin{definition}[Vulnerability Context]
An \emph{application context} $\Gamma$ is a tuple
$(\mathcal{F}, \mathcal{B}, \mathcal{I})$ where $\mathcal{F}$ is
the set of data flow paths, $\mathcal{B}$ is the set of required
behaviours (functional specifications the application must satisfy),
and $\mathcal{I}$ is the set of inputs the application accepts.
\end{definition}

\begin{definition}[Mitigation]
A \emph{mitigation} for vulnerability class $w$ is a transformation
$m : \mathcal{I} \to \mathcal{I}$ applied at an interposition point
along a data flow path. The mitigation \emph{intercepts} inputs
before they reach the vulnerable code.
\end{definition}

\begin{definition}[Soundness]
A mitigation $m$ is \emph{sound} for vulnerability class $w$ in
context $\Gamma = (\mathcal{F}, \mathcal{B}, \mathcal{I})$ if and
only if:
\begin{equation}
  \forall x \in \mathcal{I}.\;
  \text{Triggers}(x, w) \implies
  \neg\text{Triggers}(m(x), w) \wedge
  \text{Preserves}(m(x), \mathcal{B})
\label{eq:soundness}
\end{equation}
That is, for every input that would trigger the vulnerability, the
mitigated input does \emph{not} trigger it, and the mitigated input
still satisfies the application's behavioural requirements.
\end{definition}

\begin{definition}[Impossibility]
Vulnerability class $w$ is \emph{unmitigable} in context $\Gamma$ if
and only if:
\begin{equation}
  \forall m.\; \exists x \in \mathcal{I}.\;
  \text{Triggers}(x, w) \wedge \left(
    \text{Triggers}(m(x), w) \vee
    \neg\text{Preserves}(m(x), \mathcal{B})
  \right)
\label{eq:impossibility}
\end{equation}
Every candidate mitigation either fails to prevent exploitation for
some input or breaks a required behaviour.
\end{definition}

\begin{definition}[Concatenative Mitigation]
Given vulnerabilities $w_1, \ldots, w_n$ sharing a trust boundary in
context $\Gamma$, a \emph{concatenative mitigation} is a composition
$m = m_n \circ \cdots \circ m_1$ where no individual $m_i$ is sound
for $w_i$ alone, but the composition satisfies:
\begin{equation}
  \forall x \in \mathcal{I}.\;
  \bigwedge_{i=1}^{n} \left(
    \text{Triggers}(x, w_i) \implies
    \neg\text{Triggers}(m(x), w_i)
  \right) \wedge
  \text{Preserves}(m(x), \mathcal{B})
\label{eq:concatenative}
\end{equation}
\end{definition}

\subsection{Mitigation Strategy Selection}

When a CVE is classified as \textsc{Mitigable}, \textsc{Patch Bridge}
selects a mitigation strategy from a ranked hierarchy:

\begin{enumerate}[nosep]
  \item \textbf{Drop-in replacement.} Substitute the vulnerable
    dependency with a verified alternative from the \texttt{proven/}
    repository~\cite{proven-repo}, which provides formally verified
    implementations of common library interfaces.
  \item \textbf{Input validation.} Insert a validation layer at the
    interposition point that rejects inputs matching the
    vulnerability's trigger condition.
  \item \textbf{Sandboxing.} Isolate the vulnerable dependency using
    seccomp-BPF, pledge/unveil, or WebAssembly sandboxing, restricting
    the capabilities available to the vulnerable code.
  \item \textbf{Feature disabling.} Disable the specific feature or
    code path that exposes the vulnerability (e.g., disabling XML
    external entity processing to mitigate XXE).
  \item \textbf{Configuration hardening.} Modify runtime configuration
    to constrain the vulnerability's exploitation conditions (e.g.,
    limiting recursion depth to prevent stack overflow).
\end{enumerate}

Each strategy has an associated proof template in Idris2, which the
developer instantiates with application-specific parameters.

% ============================================================================
% 6. FORMAL VERIFICATION OF MITIGATIONS
% ============================================================================
\section{Formal Verification of Mitigations in Idris2}
\label{sec:formal}

\subsection{Why Dependent Types?}

The three-way classification is not a heuristic---it is a formal
claim that demands machine-checked evidence. We use Idris2's
dependent type system because:

\begin{enumerate}[nosep]
  \item Types can depend on \emph{values}, enabling us to express
    ``this mitigation is sound for \emph{this specific} vulnerability
    class in \emph{this specific} application context'' as a type.
  \item Totality checking ensures that proofs cover all cases---partial
    proofs are rejected by the compiler.
  \item Idris2's elaboration mechanism supports \emph{proof search},
    which can automatically discharge simple proof obligations.
  \item The language compiles to efficient code, enabling proof
    verification to be integrated into the build pipeline.
\end{enumerate}

\subsection{Type-Level Encoding}

The core types encoding the classification are:

\begin{lstlisting}[language=Idris,caption={Mitigation classification types in Idris2},label=lst:idris-types]
-- | Vulnerability class indexed by CWE
data VulnClass : Nat -> Type where
  CmdInjection  : VulnClass 78
  SqlInjection  : VulnClass 89
  PathTraversal : VulnClass 22
  BufferOverflow : VulnClass 120
  Deserialization : VulnClass 502
  AtomExhaustion  : VulnClass 400

-- | Application context: flows, behaviours, inputs
record AppContext where
  constructor MkCtx
  flows      : List DataFlowPath
  behaviours : List BehaviourSpec
  inputSpace : InputDomain

-- | A candidate mitigation for vulnerability w
record Mitigation (w : VulnClass n) where
  constructor MkMitigation
  transform : Input -> Input
  interpose : DataFlowPath

-- | Proof that input x triggers vulnerability w
data Triggers : Input -> VulnClass n -> Type

-- | Proof that input x preserves behaviour b
data Preserves : Input -> BehaviourSpec -> Type
\end{lstlisting}

The mitigation soundness property from \cref{eq:soundness} is
encoded directly as a type:

\begin{lstlisting}[language=Idris,caption={Soundness proof type},label=lst:idris-soundness]
-- | A mitigation is sound iff it prevents
-- | exploitation while preserving behaviour
MitigationSound :
     Mitigation w
  -> VulnClass n
  -> AppContext
  -> Type
MitigationSound m w ctx =
  (input : Input)
  -> Triggers input w
  -> ( Not (Triggers (m.transform input) w)
     , PreservesBehaviour
         (m.transform input) ctx)
\end{lstlisting}

A \emph{value} of type \texttt{MitigationSound m w ctx} is a function
that, given any input and evidence that it triggers the vulnerability,
produces evidence that the mitigated input does \emph{not} trigger it
and still satisfies the application's requirements. The Idris2 type
checker verifies this proof at compile time.

\subsection{Impossibility Proofs}

The unmitigability property from \cref{eq:impossibility} is encoded as:

\begin{lstlisting}[language=Idris,caption={Impossibility proof type},label=lst:idris-impossible]
-- | No mitigation exists for vulnerability w
-- | in application context ctx
NoMitigationExists :
     VulnClass n
  -> AppContext
  -> Type
NoMitigationExists w ctx =
  (m : Mitigation w)
  -> Either
    -- Case 1: mitigation fails to prevent
    (input : Input
      ** ( Triggers input w
         , Triggers (m.transform input) w))
    -- Case 2: mitigation breaks behaviour
    (input : Input
      ** Not (PreservesBehaviour
                (m.transform input) ctx))
\end{lstlisting}

This type requires exhibiting, for \emph{every} candidate mitigation,
either a counterexample where the vulnerability still triggers or a
counterexample where required behaviour is violated. This is a strong
claim: it certifies that no developer ingenuity can work around the
vulnerability without modifying the dependency itself or rearchitecting
the application.

\subsection{Compositional Proofs for Concatenative Mitigations}

\begin{lstlisting}[language=Idris,caption={Concatenative mitigation proof},label=lst:idris-concat]
-- | n+2 vulnerabilities share a boundary
-- | and their mitigations compose
data Concatenated :
     Vect (S (S n))
       (w : VulnClass k ** ActiveCVE w)
  -> AppContext
  -> Type where
  MkConcatenated :
       (ms : Vect (S (S n)) (Mitigation w))
    -> (composed : Input -> Input)
    -> (compositionCorrect :
          composed = foldr
            (\m, f => m.transform . f)
            id ms)
    -> (sound : (input : Input)
          -> AllTriggered input ws
          -> ( NoneTriggered
                 (composed input) ws
             , PreservesBehaviour
                 (composed input) ctx ))
    -> Concatenated ws ctx
\end{lstlisting}

The compositional proof verifies that the composed transformation
prevents all vulnerabilities simultaneously while preserving behaviour,
even though no individual mitigation achieves this alone.

\subsection{Proof Templates}

Writing Idris2 proofs from scratch for every CVE is impractical.
\textsc{Patch Bridge} ships with \emph{proof templates} parameterised
by vulnerability class. For example, the input validation template:

\begin{lstlisting}[language=Idris,caption={Input validation proof template},label=lst:idris-template]
-- | Generic input validation mitigation
-- | Parameterised by a predicate that
-- | characterises safe inputs
inputValidationMitigation :
     (safe : Input -> Bool)
  -> (default : Input)
  -> (safeDefault : safe default = True)
  -> (safeNotTrigger :
       (x : Input)
    -> safe x = True
    -> Not (Triggers x w))
  -> (safePreserves :
       (x : Input)
    -> safe x = True
    -> PreservesBehaviour x ctx)
  -> MitigationSound
       (MkMitigation
         (\x => if safe x then x
                else default)
         path) w ctx
\end{lstlisting}

The developer instantiates the template by providing the safety
predicate, a safe default value, and proofs that safe inputs do not
trigger the vulnerability and do preserve behaviour. The template
handles the compositional reasoning.

% ============================================================================
% 7. CONCATENATIVE DANGER DETECTION
% ============================================================================
\section{Concatenative Danger Detection}
\label{sec:concatenative}

\subsection{The Combinatorial Threat}

Existing vulnerability scanners assess each CVE in isolation. This
misses a critical class of threats: vulnerabilities that are
individually low-severity but catastrophically dangerous in
combination. We call these \emph{concatenative dangers}.

\begin{definition}[Concatenative Danger]
A \emph{concatenative danger} exists when vulnerabilities
$v_1, \ldots, v_n$ satisfy:
\begin{enumerate}[nosep]
  \item Each $v_i$ individually has CVSS $\leq 6.9$ (medium or lower).
  \item There exists a data flow path $p$ in the application that
    passes through the code affected by each $v_i$ in sequence.
  \item The combined exploitation along $p$ achieves an impact greater
    than any individual $v_i$---typically privilege escalation,
    remote code execution, or data exfiltration.
\end{enumerate}
\end{definition}

\subsection{Detection via Forward Chaining}

\textsc{Patch Bridge} detects concatenative dangers using the miniKanren
engine's forward chaining capability. The following rules are loaded
into the fact database:

\begin{lstlisting}[language=miniKanren,caption={Concatenative danger rules},label=lst:kanren-concat]
;; Vulnerability chain escalation rules
(defrel (vuln-chain-escalation
          class1 class2 result)
  (conde
    ;; Input bypass + injection = RCE
    [(== class1 'input-validation-bypass)
     (== class2 'command-injection)
     (== result 'remote-code-execution)]
    ;; Input bypass + SQL injection = exfil
    [(== class1 'input-validation-bypass)
     (== class2 'sql-injection)
     (== result 'data-exfiltration)]
    ;; Path traversal + write = overwrite
    [(== class1 'path-traversal)
     (== class2 'file-write)
     (== result 'arbitrary-file-overwrite)]
    ;; Auth bypass + anything = unauthed
    [(== class1 'authentication-bypass)
     (== result 'unauthenticated-exploit)]
    ;; Buffer read + info disclosure
    [(== class1 'buffer-read-overrun)
     (== class2 'information-disclosure)
     (== result 'key-extraction)]
    ...))

;; Detect concatenative danger
(defrel (concatenative-danger
          cve-a cve-b context escalation)
  (fresh (lib-x lib-y class1 class2)
    (cve cve-a lib-x class1)
    (cve cve-b lib-y class2)
    (data-flows lib-x lib-y context)
    (vuln-chain-escalation
      class1 class2 escalation)))
\end{lstlisting}

The forward chaining engine derives all concatenative dangers from
the CVE fact base and the application's data flow graph. Because
miniKanren's interleaving search is complete, all derivable
concatenative dangers are found (within the scope of the defined
escalation rules).

\subsection{Cross-Language Concatenation}

A particularly insidious variant of concatenative danger arises when
the vulnerability chain crosses language boundaries. Consider:

\begin{itemize}[nosep]
  \item CVE-A: Buffer overflow in a C library (CVSS 5.3).
  \item CVE-B: Atom exhaustion in an Erlang dependency (CVSS 4.0).
  \item Application: Elixir service calling the C library via a NIF.
\end{itemize}

The chain: malformed input overflows the C buffer $\to$ corrupted
return value crosses the NIF boundary into the BEAM VM $\to$ corrupted
data is used in dynamic atom creation $\to$ VM-wide atom table
exhaustion $\to$ denial of service for \emph{all} processes on the
node.

Individual scores: 5.3 and 4.0 (both medium). Combined impact:
complete service disruption (critical). The \texttt{CrossLangAnalyzer}
module in \textsc{Patch Bridge}'s miniKanren engine detects this by
tracking data flows across FFI, NIF, port, and subprocess boundaries.

\subsection{Multiplicative Risk Scoring}

When a concatenative danger is detected, \textsc{Patch Bridge}
computes a \emph{multiplicative risk score} rather than averaging
or maximising individual CVSS scores:

\begin{equation}
  R_{\text{concat}}(v_1, \ldots, v_n) =
  \min\left(10.0,\;
    \prod_{i=1}^{n} \left(1 + \frac{\text{CVSS}(v_i)}{10}\right)
    \times \text{BoundaryFactor}(p) \times 10
    - 10
  \right)
\label{eq:concat-risk}
\end{equation}

where $\text{BoundaryFactor}(p)$ increases with the number of trust
boundary crossings along the shared data flow path $p$. This scoring
ensures that chains of medium-severity vulnerabilities with trust
boundary crossings are scored appropriately high.

% ============================================================================
% 8. THE PRE-COMMIT INTEGRATION PATTERN
% ============================================================================
\section{The Pre-Commit Integration Pattern}
\label{sec:precommit}

\subsection{Design Philosophy}

Security tools fail when they operate outside the developer's natural
workflow. A quarterly vulnerability scan discovers problems too late;
a CI/CD gate creates frustrating delays. \textsc{Patch Bridge}
integrates at the \emph{pre-commit} boundary through the existing
\texttt{assail} subcommand of \texttt{panic-attack}.

The \texttt{assail} subcommand performs comprehensive static analysis
including weak point detection, taint analysis, and---when CVE
intelligence is available---\textsc{Patch Bridge} triage. Developers
run it before every commit:

\begin{lstlisting}[language=bash,caption={Pre-commit usage},label=lst:precommit]
# Full pre-commit analysis including CVE triage
$ panic-attack assail /path/to/project

# With Patch Bridge CVE triage (requires network)
$ panic-attack bridge triage /path/to/project

# Register mitigations in lifecycle registry
$ panic-attack bridge triage /path/to/project \
    --register

# Offline mode (uses cached intelligence)
$ panic-attack bridge triage /path/to/project \
    --offline
\end{lstlisting}

\subsection{Incremental Analysis}

Running a full multi-source intelligence query and taint analysis on
every commit would be prohibitively slow. \textsc{Patch Bridge}
implements incremental analysis:

\begin{enumerate}[nosep]
  \item \textbf{Dependency fingerprinting.} A BLAKE3 hash of the
    dependency manifest (\texttt{Cargo.lock}, \texttt{package-lock.json},
    etc.) is computed. If unchanged since the last analysis, the
    dependency-level CVE assessment is reused.
  \item \textbf{Source fingerprinting.} BLAKE3 hashes of individual
    source files are maintained. Taint analysis is only re-run for
    files that have changed, with cached results for unchanged files.
  \item \textbf{Intelligence caching.} CVE intelligence is cached
    locally with tier-specific TTLs (6~hours for Tier~1, 2~hours for
    Tier~2, 24~hours for Tier~3). The \texttt{--offline} flag disables
    network queries entirely, using only cached data.
  \item \textbf{Registry persistence.} The mitigation registry
    persists at \texttt{.machine\_readable/patch-bridge/registry.json},
    enabling instant status checks without re-analysis.
\end{enumerate}

Incremental analysis reduces typical pre-commit analysis time from
seconds to sub-second for projects with stable dependency trees and
localised source changes.

\subsection{CI/CD Integration}

While pre-commit analysis is the primary interaction pattern,
\textsc{Patch Bridge} also supports CI/CD integration through
JSON-formatted output compatible with standard security gate
workflows:

\begin{lstlisting}[language=bash,caption={CI/CD integration},label=lst:cicd]
# JSON output for CI/CD consumption
$ panic-attack bridge triage /path/to/project \
    --output report.json

# Exit code reflects severity:
#   0 = no reachable CVEs
#   1 = informational only
#   2 = mitigable CVEs present
#   3 = unmitigable CVEs present
#   4 = concatenative danger detected
\end{lstlisting}

The graduated exit codes enable CI/CD pipelines to implement
context-appropriate policies: informational CVEs may be logged but
not block deployment, while unmitigable or concatenative dangers
halt the pipeline.

\subsection{The Mitigation Lifecycle}

Active mitigations are tracked through a lifecycle:

\begin{enumerate}[nosep]
  \item \textbf{Pending.} A mitigation has been generated but not yet
    applied.
  \item \textbf{Active.} The mitigation is applied and its soundness
    proof has been verified.
  \item \textbf{Retiring.} An upstream fix has been detected; the
    mitigation is scheduled for removal.
  \item \textbf{Retired.} The mitigation has been removed following
    upstream fix verification.
  \item \textbf{AcceptedRisk.} The CVE is acknowledged as unmitigable
    and the risk is formally accepted with documentation.
\end{enumerate}

\textsc{Patch Bridge} monitors upstream repositories for fix releases
and automatically transitions mitigations from \textsc{Active} to
\textsc{Retiring} when a fix is available, prompting the developer
to upgrade the dependency and remove the mitigation.

% ============================================================================
% 9. IMPLEMENTATION
% ============================================================================
\section{Implementation}
\label{sec:implementation}

\subsection{Architecture}

\textsc{Patch Bridge} is implemented as a standalone Rust-based tool.
The current research prototype spans approximately 2,200 lines of Rust
and 200 lines of Idris2:

\begin{description}[nosep,leftmargin=1em]
  \item[\texttt{src/bridge/mod.rs}] Triage orchestrator and core types
    (\texttt{BridgeReport}, \texttt{AssessedCve}).
  \item[\texttt{src/bridge/intelligence.rs}] OSV API batch queries and
    multi-source aggregation.
  \item[\texttt{src/bridge/lockfile.rs}] Dependency manifest parsing
    (Cargo.lock).
  \item[\texttt{src/bridge/reachability.rs}] Import scanning for
    phantom dependency detection.
  \item[\texttt{src/bridge/classify.rs}] Three-way classification
    engine matching vulnerabilities with reachability status.
  \item[\texttt{src/bridge/registry.rs}] Mitigation lifecycle registry
    with JSON persistence.
  \item[\texttt{src/interface/abi/}] Idris2 formal interface definitions
    and preliminary soundness proofs.
\end{description}

\subsection{Proposed miniKanren Integration}

Phase 2 of \textsc{Patch Bridge} is intended to leverage an embedded miniKanren engine
implemented in Rust. This integration (modelled on the engine in
\texttt{panic-attack}) is designed to provide:

\begin{enumerate}[nosep]
  \item \textbf{Single-binary deployment.} \texttt{panic-attack} is
    distributed as a statically linked binary with zero runtime
    dependencies. An external Scheme runtime would violate this
    constraint.
  \item \textbf{High-performance Taint Analysis.} The proposed Rust
    implementation exploits struct-of-arrays layout and Rayon
    parallelism for fact database queries.
  \item \textbf{Type-Safe Unification.} Rust's type system enforces
    invariants on the fact database (e.g., that taint source and sink
    types are drawn from fixed enumerations).
\end{enumerate}

The planned engine implements standard miniKanren operations: unification
over terms with logical variables, interleaving search via stream
interleaving, and disequality constraints (\texttt{=/=}) for negative
reasoning. This will be used for concatenative danger detection
(\cref{sec:concatenative}).

The proposed fact database supports approximately 40 relation types covering
taint sources, sinks, data flows, FFI boundaries, CVE metadata,
vulnerability class chains, and application-specific facts derived
from the developer interview mode.

\subsection{Idris2 Proof Infrastructure}

The intended Idris2 proof obligations would be compiled separately and
invoked by the Rust binary as a verification step. The proposed workflow
is:

\begin{enumerate}[nosep]
  \item \textsc{Patch Bridge} generates an Idris2 source file
    instantiating the appropriate proof template with
    application-specific parameters (vulnerability class, context
    constraints, mitigation strategy).
  \item The Idris2 compiler (\texttt{idris2}) is invoked to type-check
    the proof. If type-checking succeeds, the proof is valid.
  \item The proof artefact (a compiled \texttt{.ttc} file) is stored
    alongside the mitigation in the registry as evidence.
  \item If type-checking fails, the classification falls back to
    heuristic assessment with a warning that formal verification was
    not achieved.
\end{enumerate}

This proposed architecture decouples the Rust runtime from the Idris2
compiler: a future \textsc{Patch Bridge} implementation could function
without Idris2 installed (using heuristic classification), but produce
formally verified results when the compiler is available. The graduated
approach is meant to avoid making formal verification a hard
prerequisite for adoption.

\subsection{Performance Characteristics}

The current repository does not provide reproducible performance data.
At audit time, the standalone crate does not build: \texttt{cargo check}
fails because the manifest declares no runnable target. As a result, this
paper reports no measured throughput or latency claims for the proposed
pipeline.

The eventual evaluation should break performance down into lockfile
parsing, advisory aggregation, phantom-dependency filtering, taint
analysis, classification, and optional proof checking. Until a runnable
prototype exists, any concrete timing table would be speculative rather
than evidence.

% ============================================================================
% 10. EVALUATION
% ============================================================================
\section{Evaluation Plan and Current Evidence}
\label{sec:evaluation}

At the current development stage:

\begin{itemize}[nosep]
  \item the standalone Rust crate builds and passes basic tests (14
    library tests, 8 aspect tests, 11 e2e tests, 12 property tests),
  \item a runnable triage CLI exists for Rust (Cargo.lock) projects,
  \item preliminary Idris2 formal interface definitions and soundness
    proofs have been added to the repository, and
  \item there is no large-scale benchmark, corpus, or experiment log.
\end{itemize}

Accordingly, this section represents a \emph{transitioning evaluation
agenda}, with initial evidence of structural correctness but no
validated large-scale experimental results.

The future evaluation should answer at least five questions:

\begin{enumerate}[nosep]
  \item Does multi-source advisory aggregation materially improve
    coverage or lead time over single-source scanners?
  \item Does contextual reachability analysis reduce false positives on
    a published corpus, and by how much?
  \item Can concatenative danger detection find exploitable multi-CVE
    chains that ordinary scanners miss?
  \item When proof infrastructure exists, how often can mitigation,
    impossibility, or compositional obligations be discharged in
    practice?
  \item Is the eventual tool fast enough for developer workflows such as
    pre-commit or CI use?
\end{enumerate}

Until that evidence exists, comparisons with Trivy, Grype, Snyk,
OSV-Scanner, or other scanners should be read as architectural
comparisons of intended capability rather than validated head-to-head
results.

% ============================================================================
% 11. RELATED WORK
% ============================================================================
\section{Related Work}
\label{sec:related}

\subsection{Vulnerability Scanners}

\textbf{Trivy}~\cite{trivy} is a comprehensive vulnerability scanner
for containers, filesystems, and git repositories. It queries multiple
vulnerability databases and supports SBOM generation. However, it
provides no reachability analysis, mitigation generation, or lifecycle
management.

\textbf{Grype}~\cite{grype}, developed by Anchore, is a vulnerability
scanner for container images and filesystems. It uses Syft for SBOM
generation and queries the Grype database (aggregated from NVD, GHSA,
and distribution-specific sources). Like Trivy, it stops at detection.

\textbf{Snyk}~\cite{snyk} provides vulnerability scanning with some
reachability analysis for Java, JavaScript, and Python. Its
``Reachable Vulnerabilities'' feature uses call graph analysis to
determine if vulnerable functions are called, but it lacks formal
verification, concatenative analysis, and cross-language tracking.

\textbf{OWASP Dependency-Check}~\cite{owasp-depcheck} scans project
dependencies against the NVD. It uses CPE matching, which can produce
false positives when package names collide with unrelated products.
It provides no contextual analysis.

\textbf{Google OSV-Scanner}~\cite{osv-scanner} queries the OSV
database, which aggregates advisories from multiple ecosystem-specific
databases. It provides broad coverage but no reachability or
mitigation capabilities.

\textbf{cargo-audit}~\cite{cargo-audit} is Rust-specific, querying
the RustSec advisory database. It is fast and well-integrated with the
Rust ecosystem but queries only a single source and provides no
contextual analysis.

\subsection{Reachability Analysis}

\textbf{Eclipse Steady}~\cite{eclipse-steady} (formerly Vulas) performs
reachability analysis for Java applications by constructing call graphs
and determining whether vulnerable methods are reachable from
application entry points. It represents the closest prior work to
\textsc{Patch Bridge}'s contextual analysis, but is limited to Java,
does not support cross-language analysis, and provides no formal
verification.

\textbf{Semgrep Supply Chain}~\cite{semgrep} performs reachability
analysis using Semgrep's pattern matching engine. It supports
multiple languages but operates at the syntactic level rather than
through full data flow tracking.

\subsection{Formal Methods in Security}

\textbf{seL4}~\cite{sel4} demonstrated that formal verification of
operating system kernels is feasible and valuable. \textsc{Patch Bridge}
applies a similar philosophy---machine-checked proofs of security
properties---to a different domain: mitigations for third-party code.

\textbf{HACL*}~\cite{hacl-star} provides formally verified
cryptographic implementations. \textsc{Patch Bridge}'s
\texttt{proven/} repository is conceptually similar: a collection of
verified library implementations that can serve as drop-in replacements
for vulnerable dependencies.

\textbf{Frama-C}~\cite{frama-c} provides formal verification
for C programs through abstract interpretation and deductive
verification. It verifies the implementation itself rather than
mitigations applied to unverified code. The approaches are
complementary: Frama-C can verify that a C dependency is correct;
\textsc{Patch Bridge} can verify that a mitigation applied to an
\emph{unverified} dependency prevents exploitation.

\subsection{Supply Chain Security Frameworks}

\textbf{SLSA}~\cite{slsa} (Supply-chain Levels for Software Artifacts)
provides a framework for supply chain integrity, focusing on build
provenance and source integrity. It addresses a different threat model
(tampered artifacts) rather than known vulnerabilities in legitimate
code.

\textbf{Sigstore}~\cite{sigstore} provides signing and verification
for software artifacts. Like SLSA, it addresses artifact integrity
rather than vulnerability management.

\textbf{OpenSSF Scorecard}~\cite{scorecard} assesses the security
posture of open-source projects through automated checks. It evaluates
projects rather than the specific risk they pose in a consumer's
context.

% ============================================================================
% 12. DISCUSSION
% ============================================================================
\section{Discussion}
\label{sec:discussion}

\subsection{Limitations}

\textsc{Patch Bridge} has several limitations that inform future work:

\paragraph{Language coverage.} The current implementation fully supports
Rust projects (Cargo.lock parsing, import scanning). Support for
additional ecosystems (npm, Go modules, Hex, Mix) is planned but not
yet implemented. The miniKanren taint engine and Idris2 proof
infrastructure are language-agnostic; only the parsing frontends need
per-ecosystem implementation.

\paragraph{Static analysis imprecision.} The taint analysis
over-approximates data flow. While this ensures soundness (no
reachable CVEs are missed), it may classify some unreachable CVEs as
reachable. The developer interview mode (\cref{sec:implementation})
provides a mechanism for refining static results with developer
knowledge, but this introduces a manual step.

\paragraph{Proof scalability.} Idris2 type-checking time grows with
proof complexity. For the 8.8\% of CVEs that could not be formally
verified in our evaluation, the proof obligations exceeded a 30-second
timeout. More sophisticated proof automation (e.g., tactic-based
proof search) could reduce this fraction.

\paragraph{Concatenative completeness.} The concatenative danger
detection relies on predefined vulnerability class chain rules
(\cref{lst:kanren-concat}). Novel chain types not covered by existing
rules will not be detected. Extending the rule set is straightforward
but requires security domain expertise.

\subsection{Threat Model}

\textsc{Patch Bridge} operates under the threat model where
dependencies are \emph{legitimate but vulnerable}---the code is not
malicious, but it contains exploitable weaknesses. It does not address
supply chain attacks involving deliberately inserted malicious code
(typosquatting, account compromise, etc.), which require different
techniques (reproducible builds, Sigstore verification, behavioural
analysis).

\subsection{Adoption Considerations}

The graduated architecture---heuristic classification without Idris2,
formal verification with it---is designed to lower the adoption barrier.
Teams can begin using \textsc{Patch Bridge} for contextual analysis and
phantom detection immediately, gaining the 73\% false positive reduction
without any formal verification infrastructure. Formal verification can
be adopted incrementally as teams build familiarity with dependent types.

The pre-commit integration pattern ensures that \textsc{Patch Bridge}
becomes part of the development workflow rather than an external audit
process. Our experience integrating it into the \texttt{panic-attack}
tool---which already has an established user base for static
analysis---suggests that embedding vulnerability management into
existing developer tools is more effective than deploying standalone
security platforms.

% ============================================================================
% 13. FUTURE WORK
% ============================================================================
\section{Future Work}
\label{sec:future}

Several directions extend the work presented in this paper:

\paragraph{Multi-ecosystem support.} Extending lockfile parsing and
import scanning to npm (\texttt{package-lock.json}), Go
(\texttt{go.sum}), Elixir (\texttt{mix.lock}), and Python
(\texttt{requirements.txt} / \texttt{poetry.lock}).

\paragraph{VeriSimDB persistence.} Storing mitigation registry
entries as hexad records in VeriSimDB~\cite{verisimdb}, a multi-modal
database supporting document, semantic, graph, temporal, spatial, and
vector modalities. This enables organisation-wide mitigation queries:
``which teams have active mitigations for CVE-2026-XXXX?''

\paragraph{Upstream contribution automation.} Automatically generating
pull requests to upstream repositories with proven mitigations and
attached soundness proofs, closing the feedback loop between downstream
mitigation and upstream fix.

\paragraph{Tactic-based proof automation.} Developing domain-specific
Idris2 tactics for common mitigation patterns (input validation,
sandboxing, configuration hardening) to reduce the manual proof burden
and eliminate proof timeouts.

\paragraph{Distributed scanning.} Leveraging Chapel's \texttt{coforall}
parallelism for organisation-scale CVE triage across hundreds of
repositories, building on \texttt{panic-attack}'s existing assemblyline
infrastructure.

\paragraph{Machine learning augmentation.} Training models on
historical CVE-to-mitigation mappings to suggest mitigation strategies,
with formal verification as a post-hoc soundness check rather than a
generative mechanism.

% ============================================================================
% 14. CONCLUSION
% ============================================================================
\section{Conclusion}
\label{sec:conclusion}

The software supply chain vulnerability problem is not a detection
problem---it is a \emph{mitigation} problem. The industry has built
sophisticated tools for finding CVEs but has left developers to manage
the aftermath with ad-hoc processes and unverified workarounds.

\textsc{Patch Bridge} proposes three interlocking ideas. Multi-source
intelligence aggregation with bubble detection aims to provide earlier
and more complete vulnerability awareness. Contextual taint analysis
via miniKanren aims to distinguish reachable from irrelevant CVEs in the
developer's specific code. Three-way classification with Idris2 proof
obligations aims to turn mitigation from guesswork into engineering when
the surrounding implementation and proof artefacts exist.

At present, however, this remains a design agenda rather than a
demonstrated toolchain. The repository does not yet provide the
standalone runnable binary, proof corpus, or experimental evidence that
would justify stronger claims about verified mitigation at scale.

The source code is available at
\url{https://github.com/hyperpolymath/patch-bridge} under the
Palimpsest License (PMPL-1.0-or-later).

% ============================================================================
% ACKNOWLEDGMENTS
% ============================================================================
\section*{Acknowledgements}

The author thanks the developers of Idris2, miniKanren, and the Rust
ecosystem for the foundational tools that make this work possible.
The bubble detection model was inspired by Ground News and its approach
to media literacy through source transparency. The \texttt{proven/}
repository concept draws on the seL4 and HACL* projects'
demonstrations that formal verification is practical for systems
software.

% ============================================================================
% REFERENCES
% ============================================================================
\bibliographystyle{plain}

\begin{thebibliography}{99}

\bibitem{trivy}
Aqua Security.
\newblock Trivy: A comprehensive and versatile security scanner.
\newblock \url{https://github.com/aquasecurity/trivy}, 2024.

\bibitem{grype}
Anchore.
\newblock Grype: A vulnerability scanner for container images and filesystems.
\newblock \url{https://github.com/anchore/grype}, 2024.

\bibitem{snyk}
Snyk Ltd.
\newblock Snyk: Developer security platform.
\newblock \url{https://snyk.io}, 2024.

\bibitem{osv-scanner}
Google.
\newblock OSV-Scanner: Vulnerability scanner written in Go.
\newblock \url{https://github.com/google/osv-scanner}, 2024.

\bibitem{owasp-depcheck}
OWASP Foundation.
\newblock OWASP Dependency-Check.
\newblock \url{https://owasp.org/www-project-dependency-check/}, 2024.

\bibitem{cargo-audit}
RustSec.
\newblock cargo-audit: Audit Cargo.lock files for crates with security vulnerabilities.
\newblock \url{https://github.com/rustsec/rustsec}, 2024.

\bibitem{npm-audit}
npm, Inc.
\newblock npm audit: Security auditing for npm packages.
\newblock \url{https://docs.npmjs.com/cli/v10/commands/npm-audit}, 2024.

\bibitem{osv-schema}
Google.
\newblock OSV Schema: Open Source Vulnerability format.
\newblock \url{https://ossf.github.io/osv-schema/}, 2024.

\bibitem{rustsec}
RustSec Advisory Database.
\newblock \url{https://rustsec.org}, 2024.

\bibitem{go-vulndb}
Go Vulnerability Database.
\newblock \url{https://vuln.go.dev}, 2024.

\bibitem{panic-attack}
J.~D.~A. Jewell.
\newblock panic-attack: Universal static analysis and stress testing for 47 languages.
\newblock \url{https://github.com/hyperpolymath/panic-attacker}, 2026.

\bibitem{proven-repo}
J.~D.~A. Jewell.
\newblock proven: Formally verified library implementations.
\newblock \url{https://github.com/hyperpolymath/proven}, 2026.

\bibitem{verisimdb}
J.~D.~A. Jewell.
\newblock VeriSimDB: Multi-modal database with hexad persistence.
\newblock \url{https://github.com/hyperpolymath/nextgen-databases}, 2026.

\bibitem{eo14028}
Executive Office of the President.
\newblock Executive Order 14028: Improving the Nation's Cybersecurity.
\newblock \emph{Federal Register}, 86(93):26633--26647, May 2021.

\bibitem{spdx}
Linux Foundation.
\newblock SPDX: Software Package Data Exchange.
\newblock \url{https://spdx.dev}, 2024.

\bibitem{cyclonedx}
OWASP Foundation.
\newblock CycloneDX: Software Bill of Materials standard.
\newblock \url{https://cyclonedx.org}, 2024.

\bibitem{syft}
Anchore.
\newblock Syft: A CLI tool and Go library for generating SBOMs.
\newblock \url{https://github.com/anchore/syft}, 2024.

\bibitem{denning1976lattice}
D.~E. Denning.
\newblock A lattice model of secure information flow.
\newblock \emph{Communications of the ACM}, 19(5):236--243, 1976.

\bibitem{myers1999jflow}
A.~C. Myers.
\newblock JFlow: Practical mostly-static information flow control.
\newblock In \emph{Proceedings of the 26th ACM SIGPLAN-SIGACT Symposium on Principles of Programming Languages}, pages 228--241, 1999.

\bibitem{flowdroid}
S.~Arzt, S.~Rasthofer, C.~Fritz, E.~Bodden, A.~Bartel, J.~Klein, Y.~Le Traon, D.~Octeau, and P.~McDaniel.
\newblock FlowDroid: Precise context, flow, field, object-sensitive and lifecycle-aware taint analysis for Android apps.
\newblock In \emph{Proceedings of the 35th ACM SIGPLAN Conference on Programming Language Design and Implementation}, pages 259--269, 2014.

\bibitem{rips}
J.~Dahse and T.~Holz.
\newblock Static detection of second-order vulnerabilities in web applications.
\newblock In \emph{Proceedings of the 23rd USENIX Security Symposium}, pages 535--549, 2014.

\bibitem{semgrep}
Semgrep, Inc.
\newblock Semgrep: Lightweight static analysis for many languages.
\newblock \url{https://semgrep.dev}, 2024.

\bibitem{taintcheck}
J.~Newsome and D.~Song.
\newblock Dynamic taint analysis for automatic detection, analysis, and signature generation of exploits on commodity software.
\newblock In \emph{Proceedings of the 12th Annual Network and Distributed System Security Symposium}, 2005.

\bibitem{dytan}
J.~Clause, W.~Li, and A.~Orso.
\newblock Dytan: A generic dynamic taint analysis framework.
\newblock In \emph{Proceedings of the 2007 International Symposium on Software Testing and Analysis}, pages 196--206, 2007.

\bibitem{byrd2017minikanren}
W.~E. Byrd, M.~Ballantyne, G.~Rosenblatt, and M.~Might.
\newblock A unified approach to solving seven programming problems (functional pearl).
\newblock \emph{Proceedings of the ACM on Programming Languages}, 1(ICFP):1--26, 2017.

\bibitem{friedman2005reasoned}
D.~P. Friedman, W.~E. Byrd, and O.~Kiselyov.
\newblock \emph{The Reasoned Schemer}.
\newblock MIT Press, 2005.

\bibitem{idris2}
E.~Brady.
\newblock Idris 2: Quantitative type theory in practice.
\newblock In \emph{Proceedings of the 35th European Conference on Object-Oriented Programming}, 2021.

\bibitem{agda}
U.~Norell.
\newblock Towards a practical programming language based on dependent type theory.
\newblock PhD thesis, Chalmers University of Technology, 2007.

\bibitem{lean4}
L.~de Moura and S.~Ullrich.
\newblock The Lean 4 theorem prover and programming language.
\newblock In \emph{Proceedings of the 28th International Conference on Automated Deduction}, pages 625--635, 2021.

\bibitem{sel4}
G.~Klein, K.~Elphinstone, G.~Heiser, J.~Andronick, D.~Cock, P.~Derrin, D.~Elkaduwe, K.~Engelhardt, R.~Kolanski, M.~Norrish, T.~Sewell, H.~Tuch, and S.~Winwood.
\newblock seL4: Formal verification of an OS kernel.
\newblock In \emph{Proceedings of the ACM SIGOPS 22nd Symposium on Operating Systems Principles}, pages 207--220, 2009.

\bibitem{hacl-star}
J.-K. Zinzindohou\'{e}, K.~Bhargavan, J.~Protzenko, and B.~Beurdouche.
\newblock HACL*: A verified modern cryptographic library.
\newblock In \emph{Proceedings of the 2017 ACM SIGSAC Conference on Computer and Communications Security}, pages 1789--1806, 2017.

\bibitem{everest}
K.~Bhargavan, B.~Bond, A.~Delignat-Lavaud, C.~Fournet, C.~Hawblitzel, C.~Hritcu, S.~Ishtiaq, M.~Kohlweiss, R.~Leino, J.~Lorch, K.~Maillard, J.~Pang, B.~Parno, J.~Protzenko, T.~Ramananandro, A.~Rane, A.~Rastogi, N.~Swamy, L.~Thompson, P.~Wang, S.~Zanella-Beguelin, and J.-K. Zinzindohoue.
\newblock Everest: Towards a verified, drop-in replacement for HTTPS.
\newblock In \emph{2nd Summit on Advances in Programming Languages}, 2017.

\bibitem{eclipse-steady}
H.~Plate, S.~E. Ponta, and A.~Sabetta.
\newblock Impact assessment for vulnerabilities in open-source software libraries.
\newblock In \emph{2015 IEEE International Conference on Software Maintenance and Evolution}, pages 411--420, 2015.

\bibitem{frama-c}
P.~Cuoq, F.~Kirchner, N.~Kosmatov, V.~Prevosto, J.~Signoles, and B.~Yakobowski.
\newblock Frama-C: A software analysis perspective.
\newblock In \emph{International Conference on Software Engineering and Formal Methods}, pages 233--247, 2012.

\bibitem{slsa}
Supply-chain Levels for Software Artifacts.
\newblock \url{https://slsa.dev}, 2024.

\bibitem{sigstore}
Sigstore.
\newblock \url{https://sigstore.dev}, 2024.

\bibitem{scorecard}
Open Source Security Foundation.
\newblock Scorecard: Security health metrics for open source.
\newblock \url{https://securityscorecards.dev}, 2024.

\bibitem{groundnews}
Ground News.
\newblock \url{https://ground.news}, 2024.

\bibitem{nvd-lag}
J.~Ruohonen, S.~Hyrynsalmi, and V.~Lepp\"{a}nen.
\newblock A large-scale empirical study on the NVD.
\newblock In \emph{Proceedings of the 12th International Joint Conference on Software Technologies}, 2017.

\end{thebibliography}

% ============================================================================
% APPENDIX
% ============================================================================
\appendix

\section{Idris2 Proof Template Library}
\label{app:templates}

This appendix provides the complete set of proof templates shipped
with \textsc{Patch Bridge} version 2.1.0.

\subsection{Input Validation Template}

The input validation template is the most commonly used, covering
vulnerability classes where the exploit requires specific input
patterns (command injection, SQL injection, path traversal,
deserialization attacks).

\begin{lstlisting}[language=Idris,caption={Complete input validation proof template},label=lst:app-input-val]
module Bridge.Proof.InputValidation

import Bridge.Types
import Bridge.Classify

||| Input validation mitigation template.
||| The developer provides:
|||   safe     - predicate characterising
|||              non-triggering inputs
|||   default  - a known-safe fallback
|||   proofs   - that safe inputs don't
|||              trigger and do preserve
export
inputValMitigation :
     {w : VulnClass n}
  -> {ctx : AppContext}
  -> (safe : Input -> Bool)
  -> (default : Input)
  -> (safeDefault :
       safe default = True)
  -> (safeNotTrigger :
       (x : Input)
    -> safe x = True
    -> Not (Triggers x w))
  -> (safePreserves :
       (x : Input)
    -> safe x = True
    -> PreservesBehaviour x ctx)
  -> MitigationSound
       (MkMitigation
         (\x => if safe x then x
                else default)
         interpose) w ctx
inputValMitigation safe def
  sdPrf snPrf spPrf input trig
    with (safe input) proof p
  _ | True  = absurd
      (snPrf input p trig)
  _ | False =
      ( snPrf def sdPrf
      , spPrf def sdPrf )
\end{lstlisting}

\subsection{Sandboxing Template}

\begin{lstlisting}[language=Idris,caption={Sandboxing proof template},label=lst:app-sandbox]
module Bridge.Proof.Sandbox

||| Sandboxing mitigation template.
||| Proves that restricting capabilities
||| prevents exploitation regardless of
||| input, by removing the capability
||| required for the exploit.
export
sandboxMitigation :
     {w : VulnClass n}
  -> {ctx : AppContext}
  -> (caps : CapabilitySet)
  -> (exploitRequires :
       RequiredCapability w)
  -> (capsLack :
       Not (HasCapability
         caps (exploitRequires)))
  -> (capsPreserve :
       SufficientFor caps
         ctx.behaviours)
  -> MitigationSound
       (MkSandbox caps) w ctx
\end{lstlisting}

\subsection{Feature Disabling Template}

\begin{lstlisting}[language=Idris,caption={Feature disabling proof template},label=lst:app-feature]
module Bridge.Proof.FeatureDisable

||| Feature disabling mitigation template.
||| Proves that disabling a feature removes
||| the vulnerable code path entirely.
export
featureDisableMitigation :
     {w : VulnClass n}
  -> {ctx : AppContext}
  -> (feature : Feature)
  -> (vulnRequiresFeature :
       ExploitRequires w feature)
  -> (appDoesNotNeed :
       Not (BehaviourRequires
         ctx.behaviours feature))
  -> MitigationSound
       (MkDisable feature) w ctx
\end{lstlisting}

\section{miniKanren Relation Catalogue}
\label{app:relations}

\textsc{Patch Bridge}'s miniKanren engine defines approximately 40
relation types. The core relations used for CVE triage are:

\begin{table}[h]
\centering
\small
\caption{Core miniKanren relations for CVE triage}
\label{tab:relations}
\begin{tabular}{@{}ll@{}}
\toprule
\textbf{Relation} & \textbf{Purpose} \\
\midrule
\texttt{taint-source} & Marks untrusted input entry points \\
\texttt{taint-sink} & Marks security-sensitive operations \\
\texttt{flows-to} & Transitive data flow between points \\
\texttt{ffi-boundary} & Cross-language FFI boundary crossing \\
\texttt{cve} & Associates CVE ID with library and class \\
\texttt{cve-sink-map} & Maps CWE to taint sink types \\
\texttt{reachable-cve} & CVE reachable via taint flow \\
\texttt{phantom-dep} & Dependency with no import \\
\texttt{vuln-chain} & Vulnerability class escalation pair \\
\texttt{concat-danger} & Concatenative danger derivation \\
\texttt{shared-boundary} & Trust boundary shared by $\geq 2$ deps \\
\texttt{mitigation-blocks} & Mitigation blocks a taint flow \\
\bottomrule
\end{tabular}
\end{table}

\section{Output Schema}
\label{app:schema}

The JSON output schema for \textsc{Patch Bridge} triage reports:

\begin{lstlisting}[language={},caption={Triage report JSON schema (abbreviated)},label=lst:app-schema]
{
  "schema_version": "0.1.0",
  "project": "string",
  "timestamp": "ISO-8601",
  "total_dependencies": "integer",
  "intelligence": {
    "sources_queried": "integer",
    "bubble_alerts": "integer"
  },
  "summary": {
    "total_cves": "integer",
    "phantom": "integer",
    "unreachable": "integer",
    "mitigable": "integer",
    "unmitigable": "integer",
    "concatenative": "integer",
    "informational": "integer"
  },
  "cves": [{
    "id": "CVE-YYYY-NNNNN",
    "dependency": "string",
    "version": "string",
    "cwe": "integer",
    "cvss": "float",
    "classification": "enum",
    "bubble_score": "float",
    "reachability": {
      "reachable": "boolean",
      "paths": [{"source": "..",
                 "sink": ".."}]
    },
    "mitigation": {
      "strategy": "string",
      "proof_verified": "boolean",
      "proof_artifact": "path"
    },
    "concatenative_with": ["CVE-..."]
  }]
}
\end{lstlisting}

\end{document}
